\section{Coverings throughout the submaximal error range}\label{sec:fullerror71}
A pairwise separation argument at radius $\delta$ requires distances
larger than $2\delta$. Since our adaptive distance is an unhalved trace
norm, this argument cannot by itself reach $\delta\ge1$. We instead
bound the number of targets that one centre can cover. The same
measurement, rather than a different binary test for each pair, will be
used in this part of the converse. Dimensions are fixed throughout. For a target family $\mathcal F$, write
$\mathcal M_N(\delta;\mathcal F)$ for its least covering cardinality
in $d_N$ with arbitrary legal memoryless centres, and add a superscript
$\Q$ when the decoded Choi blocks must be rational.

\begin{lemma}[Multiplicity of a covering centre]\label{lem:multiplicity71}
Let $J_1,\ldots,J_K$ be targets and let one common tester with at most
$N$ calls have disjoint classical events $A_1,\ldots,A_K$ such that
$\Pr_{J_i}(A_i)\ge1-\eta$. If $0\le\eta<1$ and
$0\le\delta<2(1-\eta)$, every covering of these targets by $M$ legal
centres at adaptive radius $\delta$ satisfies
\begin{equation}\label{eq:multiplicity71}
                M\ge (1-\eta-\delta/2)K.
\end{equation}
A centre need not belong to the target family. Indeed the conclusion
still holds for a centre specified by a legal $N$-slot process with
private memory, provided its distance to each assigned target controls
this same tester.
\end{lemma}
\begin{proof}
Let $Q$ be the tester's output distribution at a centre. For each target
assigned to that centre, trace contraction under the terminal
measurement gives
$Q(A_i)\ge1-\eta-\delta/2$. The division by two is required because
classical total variation is half the unhalved trace norm. Disjointness
gives at most $(1-\eta-\delta/2)^{-1}$ assigned targets. Sum this
bound over centres. No product structure, unbiasedness or independence
is required of $Q$.
\end{proof}

\begin{lemma}[One common finite-dimensional estimator]\label{lem:tomography71}
Put $b=m(dn)^2$. On $\mathcal C_{d,n,m}$ there is, for every $N\ge1$,
one nonadaptive tester and a Hermitian tuple estimator $\widehat J$
satisfying
\begin{equation}\label{eq:tomography71}
 \sup_{J\in\mathcal C_{d,n,m}}
      \E_J |\widehat J-J|_2^2\le \frac{C_0}{N},
                  \qquad C_0=2d^2b^2.
\end{equation}
The estimator need not be positive or trace preserving. The tester
and estimator are independent of the unknown target and of any later
choice of a packing set.
\end{lemma}
\begin{proof}
One normalized maximally entangled input produces the flagged Choi
state $\Omega_J=d^{-1}\bigoplus_y J_y$. Choose a Hilbert--Schmidt
orthonormal Hermitian basis $H_1,\ldots,H_b$ of the block-diagonal
space. Every $H_j$ has operator norm at most one. Thus
$(I+H_j)/2,(I-H_j)/2$ is a binary POVM, and its signed outcome has
mean $\operatorname{tr}(H_j\Omega_J)$ and variance at most one.
For $N\ge b$, use $k=\lfloor N/b\rfloor$ independent Choi inputs
for each of the $b$ measurements. There are at most $N$ calls. Let
$\widehat a_j$ be the corresponding empirical signed mean and set
$\widehat\Omega=\sum_j\widehat a_jH_j$. Orthonormality gives
\[
 \E_J\|\widehat\Omega-\Omega_J\|_F^2
   =\sum_j\operatorname{Var}(\widehat a_j)
   \le b/k\le2b^2/N.
\]
Set $\widehat J=d\widehat\Omega$, interpreted as a tuple. For $N<b$
use $\widehat J=0$ and make no calls. Then
$|J|_2^2=d^2\operatorname{tr}\Omega_J^2\le d^2\le C_0/N$.
This proves every asserted range, without a rank or eigenvalue-margin
condition. The binary POVMs are block diagonal in the public outcome,
so coherent access to different classical flags is not needed.
\end{proof}
This elementary informationally complete estimation argument is a
standard finite-dimensional tomography construction; compare the process
estimation setting of \cite{PLS66}. Here it is used only to build a common
lower-bound witness. It does not constitute a new optimal-learning
claim, and the description encoder below is still supplied its target.

\begin{theorem}[A full-range covering lower bound]\label{thm:fullpacking71}
Let $\mathcal F\subseteq\mathcal C_{d,n,m}$ contain the image of a
nonempty $a$-dimensional Euclidean ball under a bi-Lipschitz map for
tuple Frobenius distance, with $a>0$. For every fixed $0<\delta_*<2$
there is $c>0$, depending on this patch, the fixed dimensions and
$\delta_*$, such that
\begin{equation}\label{eq:fullpacking71}
       \mathcal M_N(\delta;\mathcal F)\ge cN^{a/2}
             \quad(N\ge1,\ 0<\delta\le\delta_*).
\end{equation}
Centres may be arbitrary legal memoryless instruments. The same lower
bound holds under the relaxation to legal $N$-slot centres described
in Lemma~\ref{lem:multiplicity71}.
\end{theorem}
\begin{proof}
Set $\eta=(2-\delta_*)/4>0$ and
$r_N=(C_0/(\eta N))^{1/2}$. The previous lemma and Markov's inequality
give
$\Pr_J\{|\widehat J-J|_2\le r_N\}\ge1-\eta$ for every target.
A volumetric packing of the fixed patch supplies at least
$c_1r_N^{-a}$ points at pairwise tuple distance greater than $2r_N$.
The constant can be decreased to cover coarse scales: a singleton
packing always exists. The estimator events consisting of the closed
$r_N$-balls at these points are disjoint. Lemma~\ref{lem:multiplicity71}
applies because
$1-\eta-\delta/2\ge1-\eta-\delta_*/2=\eta$.
It gives $\mathcal M_N\ge\eta c_1r_N^{-a}$, as required. In
particular one centre may cover several packing points when the error
is large; the proof bounds that multiplicity rather than asserting it
is one.
\end{proof}

\begin{corollary}[Full-range interior and preparation laws]\label{cor:fullrange71}
Fix $0<\delta_*<2$. In Theorem~\ref{thm:jointcode68}, with fixed
$d,n,m,a$ and $s=d^2(mn^2-1)>0$, both real and rational legal-centre
covering numbers satisfy
\begin{equation}\label{eq:interiorfull71}
 \log_2\mathcal M_N(\delta;a)
 =\log_2\mathcal M_N^{\Q}(\delta;a)+O(1)
 =\frac{s}{2}\log_2N+s\log_2(1/\delta)+O(1)
\end{equation}
uniformly for $N\ge1$ and $0<\delta\le\delta_*$. The constants may
depend on $\delta_*$ and the positive margin $a$.
For the preparation family in Theorem~\ref{thm:rankentropy68}, the
same extension holds with $s$ replaced by
$v=\sum_y r_y(2n-r_y)-1$, with fixed dimensions and public output-rank
bounds. Its constants require no positive eigenvalue margin. If $v=0$
the family is a singleton and needs no payload.
\end{corollary}
\begin{proof}
For $0<\delta\le1/32$ these are the inherited joint laws. The interior
ball and maximal-rank preparation chart used in their proofs are
bi-Lipschitz patches of dimensions $s$ and $v$, respectively. Theorem
\ref{thm:fullpacking71} therefore supplies the lower term
$(s/2)\log_2N+O(1)$, or $(v/2)\log_2N+O(1)$, on the remaining
range $1/32<\delta\le\delta_*$. On that range
$\log_2(1/\delta)$ is bounded above and below by constants depending
only on $\delta_*$. For the upper bound use the old code at accuracy
$1/32$, which is also accurate to $\delta$. This proves a uniform
additive remainder for both kinds of centres. If $\delta_*\le1/32$
there is no remaining range.
\end{proof}
The cap $\delta_*<2$ is substantive: at unhalved error two every
nonempty family is covered by a single legal centre. Nothing above is
uniform as $\delta_*\uparrow2$, or along a varying error sequence
converging to two. The exact zero-error endpoint is also separate from
the positive-error coding law. The older coherent/preparation theorem
retains its own $\delta\le1/32$ range; the present common tomography
argument gives the half-dimension scale, not its coherent $N^{-1}$ scale.
